%%%%%%%%%%%%%%%%%%%%%%%%%%%%%%%%%%%%%%%%%%%%%%%%%%
% FORMATING
%%%%%%%%%%%%%%%%%%%%%%%%%%%%%%%%%%%%%%%%%%%%%%%%%%
\documentclass[a4paper,12pt,oneside]{book}
\usepackage{blindtext}
\usepackage{geometry}
\usepackage{titlesec}
\usepackage{hyperref}
\usepackage{afterpage}
\usepackage{indentfirst}
\usepackage{etoolbox}
\usepackage{lmodern}
\usepackage{graphicx}
\usepackage{sectsty}
\graphicspath{ {./} }

\font\tfont=cmr12 at 14pt
\sectionfont{\fontsize{12}{15}\selectfont}

\setlength{\parindent}{2em}
\setlength{\parskip}{0.5em}
\newcommand\blankpage{%
	\null
	\thispagestyle{empty}%
	\addtocounter{page}{-1}%
	\newpage}

\hypersetup{
	colorlinks,
	citecolor=black,
	filecolor=black,
	linkcolor=black,
	urlcolor=black
}
\geometry{
	left=40mm,
	top=30mm,
	right=30mm,
	bottom=30mm,
}
\renewcommand{\chaptername}{BAB}
\renewcommand{\contentsname}{Daftar Isi}
\renewcommand{\figurename}{Gambar}
\renewcommand{\tablename}{Tabel}
\renewcommand{\listfigurename}{Daftar Gambar}
\renewcommand{\listtablename}{Daftar Tabel}

\titleformat{\chapter}[display]
{\normalfont\bfseries}
{BAB \thechapter \centering}{-2ex}
{
	\rule{\textwidth}{0pt}
	\vspace{0ex}
	\centering}
[\vspace{-5ex}\rule{\textwidth}{0pt}]

\titlespacing*{\chapter}{0pt}{-50pt}{14pt}
%%%%%%%%%%%%%%%%%%%%%%%%%%%%%%%%%%%%%%%%%%%%%%%%%%
% BEGIN BOOK
%%%%%%%%%%%%%%%%%%%%%%%%%%%%%%%%%%%%%%%%%%%%%%%%%%
\begin{document}
\title{\tfont DESAIN DAN SIMULASI\\
	PERLINDUNGAN PROPERTI INTELEKTUAL\\
	MENGGUNAKAN ALGORITME FILTER DIGITAL\\.\\	
	DESIGN AND SIMULATION\\
	OF INTELECTUAL PROPERTIES PROTECTION\\
	USING DIGITAL FILTER ALGORITHM
}
\author{
	Disusun sebagai syarat untukmemperoleh gelar Sarjana Teknik\\
	pada Program Studi S1 Sistem Komputer\\
	Universitas Telkom\\\\
	Oleh:\\\\
	\textbf{Hanjara Cahya Adhyatma}\\
	\textbf{1104131113}\\
	\includegraphics[scale=0.13]{logo}
}

\date{
	FAKULTAS TEKNIK ELEKTRO\\
	UNIVERSITAS TELKOM\\
	BANDUNG\\
	2017}
\maketitle
\tableofcontents
\listoffigures
\listoftables
\addtocounter{page}{-3}
\afterpage{\blankpage}
\chapter{Pendahuluan}
\section{Latar Belakang}

Integrated Circuit (IC) merupakan modul teknologi dasar dari perangkat elektronika tertanam modern. Dengan berkembangnya teknologi IC yang mengutamakan ukuran kecil, dan performa yang tinggi serta dengan harga yang murah membuat teknologi IC semakin diminati [1].

Dengan ukuran modul yang sangat kecil dan banyaknya komponen
pembangun, kerja sama antara desainer dilakukan untuk membangun sebuah modul VLSI sehingga setiap desainer dapat fokus mendesain salah satu fungsi yang terdapat dalam modul tersebut. Kerja sama dilakukan untuk mempermudah pembuatan desain VLSI yang memiliki tingkat kerumitan yang tinggi. Desainer juga dapat mempercepat waktu mendesain dengan menggunakan kode sumber yang sudah ada atau bekerja sama secara paralel membuat masing-masing modul yang nantinya akan digabung menjadi sebuah modul utama VLSI.

Setelah modul selesai dibuat maka modul siap untuk di-produksi. Dalam proses produksi modul perusahaan tempat desainer bekerja tidak perlu memiliki pabrik produksi modul sendiri, perusahaan dapat bekerja sama dengan mitra percetakan yang akan memproduksi modul buatan perusahaan modul tersebut. Cara kerja sama seperti ini disebut dengan Fabless Manufacturing [2]. Ketika akan memproduksi IC, perusahaan harus menyerahkan blueprint modul VLSI ke percetakan, namun blueprint tersebut tidak terjamin kerahasiaan nya serta memungkinkan plagiarisme desain oleh oknum perusahaan atau pihak ketiga yang tertarik menggunakan desain VLSI yang telah diserahkan untuk di-produksi.

Dengan memberikan rangkaian watermark sebagai pengamanan pada cetak biru VLSI siap cetak yang menandakan kepemilikan dari desainer atau perusahaan produsen modul akan melindungi dari kecurangan pihak lain yang akan mencuri desain. Sehingga kemungkinan pencurian atau plagiarisme yang menyebabkan kerugian pada perusahaan atau desainer karena desain nya dicuri atau di-plagiat berkurang

\section{Rumusan Masalah}
Berikut ini dijelaskan rumusan masalah yang dihadapi dalam proposal penelitian Intelectual Property Protection (IPP) menggunakan metode Digital \textit{Filter Algorithm using Logical Polymorph Gate Key Verification} :
\paragraph{a.} Dengan metode \textit{Fabless Manufacturing}, desain modul yang siap diproduksi diserahkan kepada perusahaan percetakan mitra sehingga mitra dapat mengetahui desain modul dari desainer yang memungkinkan desain dapat dicuri oleh oknum percetakan atau pihak ketiga yang tertarik dengan desain tersebut.
\paragraph{b.}Desain modul rawan terhadap plagiarisme karena desain elektronik sangat mudah ditiru, sehingga pengamanan desain harus dilakukan agar desain tidak mudah untuk dicuri atau di-plagiat.
\paragraph{c.}Apabila pihak ketiga mencuri desain, desainer dapat mengklaim modul tersebut dengan bukti dari pengamanan watermark yang telah tertanam dalam IC menggunakan teknik pemanggilan watermark yang hanya diketahui oleh desainer yang mendesain IC tersebut.

\section{Tujuan}
Berikut merupakan tujuan pengamanan desain modul yang siap cetak sehingga aman terhadap pencurian hak cipta :
\paragraph{a.} Merancang rangkaian pengamanan dalam sebuah chip design sebagai bukti kepemilikan desain (ownership) atau watermarking.
\paragraph{b.} Desain chip yang telah diberi rangkaian watermark akan dianalisis perubahan performa dari desain sebelum dan sesudah watermarking serta kemungkinan watermark di-modifikasi oleh pihak lain atau reverse engineering untuk digunakan kembali oleh pengguna yang tidak sah.
\paragraph{c.} Rangkaian ini akan ditanam di dalam chip yang pemanggilan informasi pemilik dari chip hanya diketahui oleh pemilik cipta.

\section{Batasan Masalah}
Dalam penelitian ini rancangan desain VLSI yang disisipkan watermark membatasi masalah serta pembahasan yang akan diteliti sebagai berikut :
\paragraph{a.} Tidak membuat modul IC VLSI spesifik, namun menggunakan yang sudah ada dan menyisipkan dengan watermark.
\paragraph{b.} Menyisipkan rangkaian dengan data watermark dan tidak membahas detail data dari pemilik cipta.
\paragraph{c.} Watermarking yang dilakukan untuk satu chip IC dan tidak mewatermark masing-masing modul yang ter-integrasi dalam chip IC.

\section{Hipotesis}
\blindtext[1]
\begin{table}
	\caption{table A}
\end{table}
\begin{table}
	\caption{table B}
\end{table}
\afterpage{\blankpage}
%-----------------------------------------------------------------------------%
\chapter{\babDua}
%-----------------------------------------------------------------------------%

%-----------------------------------------------------------------------------%
\section{Pekerjaan Sebelumnya dan Keterkaitan}
%-----------------------------------------------------------------------------%
Secara garis besar teknik Intelectual Property Protection (IPP)
watermarking dapat diklasifikasikan menjadi 2 kelas yaitu Dynamic
Watermarking dan Static Watermarking. Dynamic Watermarking merupakan
watermark yang tidak dapat terdeteksi kecuali dengan menjalankan IP yang telah
di-watermark untuk mendeteksi sinyal yang dihasilkan, seperti digital signal
processing (DSP), atau finite state mechine (FSM) watermarking. Static
Watermarking merupakan watermark yang mengacu pada properti dari sebuah
desain, dan hanya bisa terdeteksi dengan cara statis yang berbeda, seperti jalur dan
penempatan watermarking [7].

Salah satu pengamanan lain adalah mengonversi fail simulasi dari fail.
RTL source code yang memungkinkan tidak mudah untuk di-reverse-engineering
oleh pihak ketiga, sehingga model tidak dapat dirubah dan digunakan kembali
dengan keperluan lain oleh pihak ketiga dan pengguna yang tidak bertanggung
jawab.[8][9]

Namun cara tersebut hanya melindungi dari sisi softwere yang melindungi
IP agar tidak di-salah-gunakan oleh pengguna pihak ketiga. Untuk pengamanan IP
yang digunakan dalam sharing project dan reusable project dapat digunakan
dengan pengamanan Digital Signal Processing cell yang memungkinkan integrasi
dalam sistem.

Dalam penelitian kali ini akan melakukan kombinasi dari proteksi IP
polimorph gate dengan algoritme filter digital. Menggunakan gabungan dari dua
teknik ini akan memberikan tambahan keamanan pada proteksi IP yang
kemungkinan tingkat over write watermark lebih kecil. Oleh karena itu dalam
penelitian ini mengajukan sebuah gabungan metode yang sudah ada untuk
meningkatkan kemampuan pengamanan dalam sebuah modul VLSI yang sudah
ada. Dengan menggabungkan polygate sebagai kunci kombinasi untuk
mengaktifkan modul filter digital yang akan digunakan sebagai watermark.

\section{Perancangan dan Implementasi Algoritme DSP untuk IPP}

Melakukan analisis terhadap masalah yang dikaji kemudian akan
dilakukan rancangan Intelectual Property Protection (IPP) dengan algoritme Filter
Digital yang dibangun meliputi rangkaian uji. Dari desain modul VLSI yang telah
ada akan diuji coba kan performa sebelum diberi watermark.

Dengan memberikan rangkaian watermark sebagai pengamanan pada
blueprint VLSI siap cetak yang menandakan kepemilikan dari desainer atau
perusahaan produsen modul akan melindungi dari kecurangan pihak lain yang
akan mencuri desain tersebut. Sehingga kemungkinan pencurian atau plagiarisme
berkurang yang menyebabkan kerugian pada perusahaan atau desainer karena
desain nya dicuri atau diplagiat.

Desain akan dirancang dengan kombinasi Low Pass Filter, High Pass
Filter, Band Pass Filter, dan Band Reject Filter. Kombinasi ini akan ditentukan
dan diaktifkan oleh polygate sebagai kunci pengaktifan kombinasi Filter digital.
Setelah Filter digital aktif maka kombinasi data akan melewati kombinasi filter
yang diaktifkan dari kombinasi polygate. Kemudian data hasil kombinasi proses
ini akan membentuk pola khusus yang menjadi data watermark dari desainer yang
mencirikan identitas desainer. Setelah diberikan watermark maka modul akan
diuji coba kan kembali performa nya. Bila terjadi penurunan performa maka akan
dilakukan perbaikan algoritma kemudian dilakukan diuji kembali performa nya.
Hingga didapat performa yang paling baik dari beberapa uji coba yang akan
dilakukan.

\afterpage{\blankpage}
\chapter{Desain dan Simulasi}
\blindtext[2]

\section{Studi Literatur}
\blindtext[2]
\section{Analisis}
\blindtext[3]
\section{Perancangan}
\blindtext[6]
\section{Pengujian}
\blindtext[3]
\section{Keluaran}
\blindtext[2]
\begin{figure}
	\caption{fig A}
\end{figure}
\begin{figure}
	\caption{fig B}
\end{figure}

\afterpage{\blankpage}
\chapter{Pengujian dan Analisis}
\blindtext[2]
\afterpage{\blankpage}
%-----------------------------------------------------------------------------%
\chapter{\babLima}
%-----------------------------------------------------------------------------%
\todo{Tambahkan kata-kata pengantar bab 5 disini.}


%-----------------------------------------------------------------------------%
\section{Mengubah Tampilan Teks}
%-----------------------------------------------------------------------------%
Beberapa perintah yang dapat digunakan untuk mengubah tampilan adalah: 
\begin{itemize}
	\item \bslash f \\
		Merupakan alias untuk perintah \bslash textit, contoh 
		\f{contoh hasil tulisan}.
	\item \bslash bi \\
		\bi{Contoh hasil tulisan}.
	\item \bslash bo \\
		\bo{Contoh hasil tulisan}.
	\item \bslash m \\
		\m{Contoh hasil tulisan}.
	\item \bslash mc \\
		\mc{Contoh hasil tulisan}.
	\item \bslash code \\ 
		\code{Contoh hasil tulisan}.
\end{itemize}


%-----------------------------------------------------------------------------%
\section{Memberikan Catatan}
%-----------------------------------------------------------------------------%
Ada dua perintah untuk memberikan catatan penulisan dalam dokumen yang Anda 
kerjakan, yaitu: 
\begin{itemize}
	\item \bslash todo \\
		Contoh: \\ \todo{Contoh bentuk todo.}
	\item \bslash todoCite \\ 
		Contoh: \todoCite
\end{itemize}


%-----------------------------------------------------------------------------%
\section{Menambah Isi Daftar Isi}
%-----------------------------------------------------------------------------%
Terkadang ada kebutuhan untuk memasukan kata-kata tertentu kedalam Daftar Isi. 
Perintah \bslash addChapter dapat digunakan untuk judul bab dalam Daftar isi. 
Contohnya dapat dilihat pada berkas thesis.tex.


%-----------------------------------------------------------------------------%
\section{Memasukan PDF}
%-----------------------------------------------------------------------------%
Untuk memasukan PDF dapat menggunakan perintah \bslash inpdf yang menerima satu 
buah argumen. Argumen ini berisi nama berkas yang akan digabungkan dalam 
laporan. PDF yang dimasukan degnan cara ini akan memiliki header dan footer 
seperti pada halaman lainnya. 

\inpdf{include}

Cara lain untuk memasukan PDF adalah dengan menggunakan perintah \bslash putpdf 
dengan satu argumen yang berisi nama berkas pdf. Berbeda dengan perintah 
sebelumnya, PDF yang dimasukan dengan cara ini tidak akan memiliki footer atau 
header seperti pada halaman lainnya. 

\putpdf{include}


%-----------------------------------------------------------------------------%
\section{Membuat Perintah Baru}
%-----------------------------------------------------------------------------%
Ada dua perintah yang dapat digunakan untuk membuat perintah baru, yaitu: 
\begin{itemize}
	\item \bslash Var \\
		Digunakan untuk membuat perintah baru, namun setiap kata yang diberikan
		akan diproses dahulu menjadi huruf kapital. 
		Contoh jika perintahnya adalah \bslash Var\{adalah\} makan ketika 
		perintah \bslash Var dipanggil, yang akan muncul adalah ADALAH. 
	\item \bslash var \\
		Digunakan untuk membuat perintah atau baru. 
\end{itemize}


\afterpage{\blankpage}
\end{document}
